\section{Actual power-cover actions on the original adjoint-defect
receiver}\label{actual-power-cover-actions-on-the-original-adjoint-defect-receiver}

Complete independent local derivation. Proof locators
\textbf{DCA0--DCA12}.

\subsection{DCA0. The object, provenance, and
category}\label{dca0.-the-object-provenance-and-category}

This calculation continues the particular receiver of RGR11. The
supporting datum remains \(\tau\langle Z_1;\mathrm{no}\ Z_2\rangle\).
All coefficient spaces, complex coordinates, positive integers, and
operators below occur after the original arithmetic reconstruction. The
two original branches retain their separate counters. No numerical
coordinate or operation is assigned to the support.

The inputs read for this calculation are
\href{https://github.com/KokunoYumeto/zeta-function-research-reader/blob/main/workbenches/splitzero-tandem/continuations/20260924-cohomology-weight-and-character-lifts/gct-weight-control/proofs/CC_RESIDUE_GYSIN_RECEIVER_INDEPENDENT.md}{\nolinkurl{CC_RESIDUE_GYSIN_RECEIVER_INDEPENDENT.md}},
RGR1, RGR9--RGR11;
\href{https://github.com/KokunoYumeto/zeta-function-research-reader/blob/main/workbenches/splitzero-tandem/continuations/20260924-cohomology-weight-and-character-lifts/gct-weight-control/proofs/CC_CONTINUOUS_SUPPORT_DUAL_SHEAF.md}{\nolinkurl{CC_CONTINUOUS_SUPPORT_DUAL_SHEAF.md}},
CSD3--CSD6 and CSD9--CSD11;
\href{https://github.com/KokunoYumeto/zeta-function-research-reader/blob/main/workbenches/splitzero-tandem/continuations/20260924-cohomology-weight-and-character-lifts/gct-weight-control/proofs/CC_LOCAL_NEARBY_CYCLES_AND_SPECIALIZATION.md}{\nolinkurl{CC_LOCAL_NEARBY_CYCLES_AND_SPECIALIZATION.md}},
LNC5--LNC6; the retained GTR1--GTR4 proof in
\href{CC_RAMIFIED_TRACE_AND_RESIDUE_ACTION.md}{\nolinkurl{CC_RAMIFIED_TRACE_AND_RESIDUE_ACTION.md}};
\href{CC_ADJOINT_DEFECT_INVARIANT_CYCLE_CROSS.md}{\nolinkurl{CC_ADJOINT_DEFECT_INVARIANT_CYCLE_CROSS.md}},
ADC1--ADC6; and the continuous-dual comparison in
\href{CC_ADJOINT_DEFECT_CONTINUOUS_DUAL.md}{\nolinkurl{CC_ADJOINT_DEFECT_CONTINUOUS_DUAL.md}},
ACD1--ACD4. These are exact programme dependencies, not new
primary-source readings. The human-source identity for the original
coefficient sheaf remains Alain Connes and Caterina Consani,
\emph{Schemes over} \(\mathbb F_1\) \emph{and zeta functions},
\href{https://arxiv.org/abs/0903.2024v3}{arXiv:0903.2024v3, §5}, with
the receiver comparison proved in CSP. The cover actions on the new
receiver are proved below.

Derived statements use sheaves of complex vector spaces. Every
coefficient operator, finite-coordinate homotopy, and specified current
map is continuous. The continuous dual is never replaced by an algebraic
dual. The proof does not assume a finite-rank perverse classification, a
strong-dual open mapping theorem, or pullback of arbitrary distributions
through a branch point. The notation \(K[1]\) puts a degree-zero
coefficient \(K\) in degree \(-1\).

\subsection{DCA1. The actual arithmetic map and its two
covariances}\label{dca1.-the-actual-arithmetic-map-and-its-two-covariances}

Retain \[
S=\{f\in\mathcal S(\mathbb R;\mathbb C):f(-v)=f(v),\ f(0)=0,\ \int_{\mathbb R}f(v)\,dv=0\},
\] \[
A=\{b\in C^\infty(\mathbb R_{>0}):
\sup_{u>0}(u^N+u^{-N})|(u\partial_u)^j b(u)|<\infty
\text{ for every }N,j\geq0\},
\] \[
\Sigma f(u)=2\sum_{k\geq1}f(ku),\qquad
J=\Sigma S,\quad Q=A/J,\quad q:A\to Q,\quad
B=A'_\beta,\quad j=q':Q'_\beta\hookrightarrow B.
\tag{DCA1.1}
\] The accepted original-image theorem gives closed \(J\) and the
continuous inverse of \(\Sigma:S\to J\). The original Mellin map and
divisor are \[
\Theta b(s)=\frac12\int_0^\infty b(u)u^s\frac{du}{u},\qquad
\Theta\Sigma f(s)=\zeta(s)\int_0^\infty f(v)v^s\frac{dv}{v}
\quad(\Re s>1),
\] \[
f_0(v)=\frac\pi2v^2(2\pi v^2-3)e^{-\pi v^2},\qquad
\Theta\Sigma f_0(s)=\frac{s(s-1)}8\pi^{-s/2}\Gamma(s/2)\zeta(s).
\tag{DCA1.2}
\] The latter entire function has values \(1/8\) at both \(0,1\), value
\(\pi/24\) at both \(-1,2\), and \[
\Theta\Sigma f_0(-2k)
=\frac{k(2k+1)(-1)^k\pi^k}{2\,k!}\zeta'(-2k)
\quad(k\geq1).
\tag{DCA1.3}
\] At an actual nontrivial zero \(\rho\) of multiplicity \(m_\rho\), its
full local expression is \[
t^{m_\rho}\frac{(\rho+t)(\rho+t-1)}8
\pi^{-(\rho+t)/2}\Gamma((\rho+t)/2)
\frac{\zeta(\rho+t)}{t^{m_\rho}}.
\tag{DCA1.4}
\] The local factor and every derivative remain part of the original
comparison. The quotient \(Q\) retains all vanishing-jet conditions
through order \(m_\rho-1\). It is not replaced by its value quotient in
the argument below. The original function is still \[
\zeta(s)=1+\sum_{k\geq2}k^{-s}=\prod_p(1-p^{-s})^{-1},\qquad
-\frac{\zeta'(s)}{\zeta(s)}
=\sum_p\sum_{k\geq1}(\log p)p^{-ks}\quad(\Re s>1).
\tag{DCA1.5}
\]

Let \(\mathscr Z\) be the actual distinct nontrivial zeros, put
\(\rho^\#=1-\overline\rho\), and retain \[
H=\ell^2(\mathscr Z,m),\qquad
\langle x,y\rangle_+=\sum_{\rho\in\mathscr Z}m_\rho x_\rho\overline{y_\rho},\qquad
(\mathsf Jy)_\rho=y_{\rho^\#}.
\tag{DCA1.6}
\] The actual value map \(E:Q\to H\) has
\((EF)_\rho=(\Theta_QF)(\rho)\), is continuous, and has dense image by
the retained full-jet interpolation theorem. The map \[
A_H(y)(F)=\langle EF,\mathsf Jy\rangle_+
\tag{DCA1.7}
\] is continuous into \(Q'_\beta\), anti-linear, and injective. For a
bounded subset \(K\subset Q\), its strong seminorm is at most
\(\sup_{F\in K}\|EF\|_+\,\|y\|_+\); density of \(E(Q)\) proves
injectivity.

For a fixed recovered parameter \(r>1\), use the actual operator \[
(T_ry)_\rho=r^\rho y_\rho,\quad U_r=rT_{1/r},\quad
D_r=T_r^*-U_r,
\] \[
(D_ry)_\rho=d_r(\rho)y_\rho,\qquad
d_r(\rho)=e^{-i\Im\rho\log r}
\bigl(r^{\Re\rho}-r^{1-\Re\rho}\bigr).
\tag{DCA1.8}
\] The known open strip makes these bounded. Set \[
C=\overline H,\qquad
\sigma_r(\overline y)=A_H(D_ry),\qquad
a=a_r=j\sigma_r:C\longrightarrow B.
\tag{DCA1.9}
\] This is exactly RGR11.3. The parameter \(r\) is fixed; the cover
degree \(n\) below is a separate positive integer.

On \(A\), write \(T_t b(u)=b(u/t)\). It preserves \(J\), descends to
\(T_t^Q\), and its value action is the displayed diagonal \(T_t\).
Directly on the actual zero coordinates, \[
T_t^*=\mathsf J U_t\mathsf J,\qquad
U_t^*=\mathsf J T_t\mathsf J.
\tag{DCA1.10}
\] For example
\((\mathsf J U_t\mathsf J y)_\rho=t^{1-\rho^\#}y_\rho=t^{\overline\rho}y_\rho\),
the positive Hilbert adjoint of multiplication by \(t^\rho\). Hence \[
(T_t^Q)'A_H(y)(F)
=\langle T_tEF,\mathsf Jy\rangle_+
=\langle EF,T_t^*\mathsf Jy\rangle_+
=A_H(U_ty)(F).
\tag{DCA1.11}
\] Exchanging \(T_t,U_t\) proves the other covariance. Since the
original diagonal \(D_r\) commutes with both operators, for \[
R=T_n',\quad R_-=T_{1/n}',\qquad
S_n=\overline U_n:C\to C,\quad L_n=\overline T_n:C\to C
\tag{DCA1.12}
\] one obtains the complete identities \[
\boxed{aS_n=Ra,\qquad aL_n=nR_-a,\qquad
RR_-=R_-R=1_B,\qquad S_nL_n=L_nS_n=n1_C.}
\tag{DCA1.13}
\] Here conjugation specifies the linear operators on the conjugate
vector space, rather than erasing the anti-linearity of (DCA1.7). All
transposes are strong-continuous because their primal maps take bounded
sets to bounded sets.

\subsection{DCA2. The current cone and the positive support
coordinate}\label{dca2.-the-current-cone-and-the-positive-support-coordinate}

Fix a positively oriented pole disk \(D\), with coordinate \(z_p=z\) at
zero or \(z_p=1/z\) at infinity, and \(i:\{p\}\hookrightarrow D\). The
coefficient current complex is \[
\mathscr T_A^k(U)=(\mathcal E_c^{-k}(U;A))',\qquad
dT(\phi)=(-1)^{k+1}T(d\phi),\qquad -2\leq k\leq0.
\tag{DCA2.1}
\] Its positive representatives are \[
c_b(\omega)=b\left(\int\omega\right),\qquad
h_pb(\eta)=b\left(\int\frac{d\arg z_p}{2\pi}\wedge\eta\right),\qquad
\delta_pb(f)=b(f(p)).
\tag{DCA2.2}
\] The angular coefficient is locally integrable, and Stokes on a
punctured disk gives \(dh_pb=\delta_pb\). The compact-test contraction
in CSD2 identifies \(\mathscr T_A\simeq\underline B[2]\), with the
positive orientation above.

The actual receiver is \[
\mathscr E=\mathscr E_r=\operatorname{Cone}(-\delta_pa)[-1],
\] \[
\mathscr E^{-1}=\mathscr T_A^{-2},\quad
\mathscr E^0=\mathscr T_A^{-1}\oplus i_*C,\quad
\mathscr E^1=\mathscr T_A^0,
\] \[
d^{-1}b=(-d_Tb,0),\qquad d^0(b,x)=-d_Tb+\delta_pax.
\tag{DCA2.3}
\] The quotient coordinate is \(x\), so its connecting arrow is
\(+\mathrm{Gys}_Ba\): \[
\underline B[1]\longrightarrow\mathscr E\longrightarrow i_*C
\xrightarrow{+\mathrm{Gys}_Ba}\underline B[2].
\tag{DCA2.4}
\] The closed stalk representatives \(b\mapsto c_b\) and
\(x\mapsto(h_pax,x)\) give \[
i^*\mathscr E\simeq B[1]\oplus C,
\qquad R\Gamma(D^*,\mathscr E)\simeq B[1]\oplus B[0],
\qquad\mathrm{res}=(1_B,a).
\tag{DCA2.5}
\] The second punctured coefficient is measured by the positive angular
current. In the restriction fibre convention
\(d(s,t)=(ds,\mathrm{res}(s)-dt)\), cancelling its identity pair gives
\[
i^!\mathscr E\simeq[C\xrightarrow{+a}B]
\quad\text{in degrees }0,1.
\tag{DCA2.6}
\] The sign is also seen directly: \(d(h_pb,0)=(-\delta_pb,h_pb)\), so
the positive Dirac class and the positive angular boundary coordinate
coincide. For shifted nearby and vanishing functors \(\Psi=R\psi[-1]\),
\(\Phi=R\phi[-1]\), \[
\Psi\mathscr E=B,\qquad\Phi\mathscr E=C,\qquad
\mathrm{can}=0,\quad\mathrm{var}=+a.
\tag{DCA2.7}
\] These statements are consequences of the current and restriction
cones, not applications of a finite-dimensional classification.

\subsection{DCA3. The full cover and a continuous cut
model}\label{dca3.-the-full-cover-and-a-continuous-cut-model}

Let \(b_n:D_s\to D\), \(z_p\mapsto z_p^n\), with the source disk chosen
as the inverse image of the target disk. This is a finite proper map.
Its direct-image stalk is the finite direct sum of the source stalks
above that point: small disjoint neighborhoods of the finite fibre
contain the inverse image of some target neighborhood, by properness.
This stalk formula proves exactness of \(b_{n*}\) on vector-space
sheaves without a finite-rank assumption. At the pole there is one
inverse-image point.

The inverse image of the target universal cover has exactly \(n\)
components. Label them \(0,\ldots,n-1\) in positive cyclic order and put
\[
V_n=B^n,\quad
\mathsf t(v_0,\ldots,v_{n-1})=(v_{n-1},v_0,\ldots,v_{n-2}),
\] \[
\Delta b=(b,\ldots,b),\qquad
\Sigma_n(v)=\sum_{j=0}^{n-1}v_j,\qquad
e_0b=(b,0,\ldots,0).
\tag{DCA3.1}
\] Thus \(\Psi(b_{n*}\mathscr E)=V_n\), with monodromy \(\mathsf t\). In
this section \(\Sigma_n\) is a finite sheet sum, different from the
original summation source \(\Sigma\) in (DCA1.1).

The circle cochain model of the punctured direct image is \[
N_n=[V_n\xrightarrow{\mathsf t-1}V_n]
\quad\text{in degrees }-1,0.
\tag{DCA3.2}
\] It follows by cutting a positive circle once, recording its constant
value and the difference of its values across the cut. Constant sections
are \(\Delta B\), and the positive loop class is measured by
\(\Sigma_n\). Explicitly \[
\ker(\mathsf t-1)=\Delta B,\quad
\operatorname{im}(\mathsf t-1)=\ker\Sigma_n,\quad
V_n/(\mathsf t-1)V_n\xrightarrow{\ [v]\mapsto\Sigma_nv\ }B.
\tag{DCA3.3}
\] For the middle equality, if \(\sum_jw_j=0\), set \(v_0=0\) and
\(v_j=-\sum_{k=1}^jw_k\) for \(1\leq j\leq n-1\). Then
\(v_{j-1}-v_j=w_j\) and \(v_{n-1}-v_0=w_0\), proving
\((\mathsf t-1)v=w\). This solution is a finite continuous operator.

The stalk model is \(S^{\mathrm{st}}=B[1]\oplus C\). In the cut model
its restriction has the representative \[
r_n^{-1}=\Delta:B\to V_n,\qquad
r_n^0=e_0a:C\to V_n.
\tag{DCA3.4}
\] The first formula is restriction of a constant current to every
sheet. The second has positive loop coordinate \(\Sigma_ne_0ax=ax\),
exactly the original angular restriction. Any representative produced
from a different cut differs by an element of \(\ker\Sigma_n\); the
explicit solution in (DCA3.3) gives its continuous chain homotopy to
(DCA3.4). The cut therefore changes the representative, not the original
angular map.

For clarity about what these data construct, if
\(j_D:D^*\hookrightarrow D\), a complex with punctured restriction
\(L\), stalk \(S^{\mathrm{st}}\), and restriction arrow
\(r:S^{\mathrm{st}}\to i^*Rj_{D*}L\) is represented by the homotopy
fibre of \[
Rj_{D*}L\oplus i_*S^{\mathrm{st}}
\longrightarrow i_*i^*Rj_{D*}L,\qquad (u,s)\longmapsto u-rs.
\tag{DCA3.5}
\] On the puncture this fibre is \(L\); at the pole it is the fibre of
\((1,-r):N\oplus S^{\mathrm{st}}\to N\), which retracts to
\(S^{\mathrm{st}}\). The natural comparison with the original complex is
therefore a stalkwise quasi-isomorphism. A map of punctured complexes, a
map of stalk complexes, and an explicit homotopy between their
restriction composites consequently give an actual derived sheaf
morphism via this fibre formula. No classification theorem is needed for
that construction.

The associated restriction fibre is \[
B\xrightarrow{b\mapsto(0,\Delta b)}C\oplus V_n
\xrightarrow{(x,v)\mapsto e_0ax-(\mathsf t-1)v}V_n
\quad(\deg-1,0,1).
\tag{DCA3.6}
\] The first arrow is a split continuous embedding; for example
projection to the zeroth component splits \(\Delta\). Quotienting this
identity pair retains \[
\Phi(b_{n*}\mathscr E)=C\oplus D_n,\qquad
D_n=V_n/\Delta B,
\] \[
\boxed{\mathcal V_n(x,[v])=e_0ax-(\mathsf t-1)v.}
\tag{DCA3.7}
\] This is well defined because \((\mathsf t-1)\Delta=0\). The canonical
arrow is \(\mathcal C_n(v)=(0,[v])\). Our positive support-cone
convention gives \[
\mathcal V_n\mathcal C_n=1-\mathsf t,
\qquad
\mathsf t_\Phi(x,[v])=(x,[\mathsf t v-e_0ax]),\qquad
\mathcal C_n\mathcal V_n=1-\mathsf t_\Phi.
\tag{DCA3.8}
\] These formulas fix both signs. Iterating the middle formula gives \[
\mathsf t_\Phi^k(x,[v])
=\left(x,\left[\mathsf t^kv-\sum_{\ell=0}^{k-1}\mathsf t^\ell e_0ax\right]\right),
\qquad\mathsf t_\Phi^n=1,
\tag{DCA3.9}
\] since \(\sum_{\ell=0}^{n-1}\mathsf t^\ell e_0ax=\Delta ax\). This
retains the full vanishing object and its monodromy, including the cover
contribution \(D_n\).

\subsection{DCA4. The genuine proper-current
map}\label{dca4.-the-genuine-proper-current-map}

On target tests define \[
(F_nT)(\phi)=T(T_n b_n^*\phi).
\tag{DCA4.1}
\] Properness preserves compact support, and smooth pullback followed by
the continuous coefficient operator is continuous on the fixed-support
test spaces. Transposition is therefore continuous for both current dual
topologies. It commutes with the Hom differential. Its representative
factors are \[
F_nc_b=c_{nRb},\qquad
F_nh_pb=h_pRb,\qquad F_n\delta_pb=\delta_pRb.
\tag{DCA4.2}
\] For the first equality, the integral of the pulled-back compact
two-form is \(n\) times its integral. On the puncture each inverse
branch contributes \(1/n\) of the positive angular form, so the sum of
all \(n\) branch contributions is that same form. Equality extends as a
locally integrable current through the pole. Evaluation at the single
inverse-image pole proves the last equality. Thus the three displayed
factors have separate geometric proofs.

By (DCA1.13), \[
\boxed{\mathsf F_n:b_{n*}\mathscr E\to\mathscr E,\qquad
(b,x)\longmapsto(F_nb,S_nx)}
\tag{DCA4.3}
\] is a current chain map: its attaching term satisfies
\(F_n\delta_pax=\delta_pRax=\delta_paS_nx\). Its stalk map, in the
actual representatives of (DCA2.5), is \[
i^*\mathsf F_n=(nR,S_n):B[1]\oplus C\longrightarrow B[1]\oplus C.
\tag{DCA4.4}
\] The source nearby space is all of \(V_n\), and the map is \[
\boxed{\mathsf F_{n,\Psi}=R\Sigma_n:V_n\to B.}
\tag{DCA4.5}
\] A single inverse branch is an orientation-preserving local
isomorphism and contributes \(R\), with no additional local degree;
summing the branches proves the formula. It intertwines \(\mathsf t\)
with identity.

On \(N_n\to N_1\) the two cochain maps are both \(R\Sigma_n\). They
commute with restriction, since \[
R\Sigma_n\Delta=nR,\qquad
R\Sigma_ne_0a=Ra=aS_n.
\tag{DCA4.6}
\] Thus the vanishing and costalk maps are exactly \[
\boxed{\mathsf F_{n,\Phi}(x,[v])=S_nx,\qquad
i^!\mathsf F_n=(S_n,R):[C\xrightarrow{a}B]\to[C\xrightarrow{a}B].}
\tag{DCA4.7}
\] The extra summand \(D_n\) in the literal vanishing source maps to
zero; it remains part of that source. The variation square is the
identity \[
R\Sigma_n\mathcal V_n(x,[v])=Rax=aS_nx.
\tag{DCA4.8}
\]

\subsection{DCA5. Invariants, coinvariants, and actual Gysin
adjunction}\label{dca5.-invariants-coinvariants-and-actual-gysin-adjunction}

On nearby invariants, (DCA4.5) restricts to \[
\Delta B\to B,\qquad\Delta b\longmapsto nRb.
\tag{DCA5.1}
\] On coinvariants, using the exact positive identification (DCA3.3), it
is \[
B\to B,\qquad b\longmapsto Rb.
\tag{DCA5.2}
\] The diagonal is not the entire nearby space; its composite with the
coinvariant quotient is multiplication by \(n\). These two functors
therefore see the two different factors of the same genuine map.

The Gysin square is \[
\begin{array}{ccc}
i_*C&\xrightarrow{\ b_{n*}(\mathrm{Gys}_Ba)\ }&b_{n*}\underline B[2]\\
\downarrow S_n&&\downarrow F_n\\
i_*C&\xrightarrow{\ \mathrm{Gys}_Ba\ }&\underline B[2].
\end{array}
\tag{DCA5.3}
\] At the level of currents this is \(F_n\delta_pa=\delta_paS_n\). Under
the oriented support adjunction \(i^!\underline B[2]=B\), its
coefficient equality is \[
Ra=aS_n.
\tag{DCA5.4}
\] The ordinary stalk factor \(nR\) belongs instead to (DCA5.1).
Replacing (DCA5.4) by \(nRa=aS_n\) would change the functor and would be
a false Gysin square. The distinction is proved by the original Dirac
and angular currents, not assigned by a weight label.

\subsection{DCA6. Construction of the transfer-adjoint direction without
branch-point current
pullback}\label{dca6.-construction-of-the-transfer-adjoint-direction-without-branch-point-current-pullback}

There is a derived sheaf map \[
\mathsf L_n:\mathscr E\longrightarrow b_{n*}\mathscr E
\tag{DCA6.1}
\] whose generic coefficient is \(R_-\), whose source coefficient on
\(C\) is \(L_n\), and whose positive supported coefficient is \(nR_-\).
We construct its complete gluing data.

Use the constant-sheaf unit with coefficient \(R_-\). On the puncture
its map is \(\Delta R_-\) into the full \(n\)-sheet coefficient. On the
pole stalk it is \(R_-\). On positive support it is \(nR_-\): a positive
boundary circle is mapped with degree \(n\), or equivalently
\(b_n^*(d\arg z_p/(2\pi))=n\,d\arg z_p/(2\pi)\). Its Gysin square is
exactly \[
nR_-a=aL_n,
\tag{DCA6.2}
\] already proved for the actual arithmetic map in (DCA1.13).

In the explicit cut models use \[
f_{\mathrm{st}}=(R_-,L_n),\qquad
f_N^{-1}=\Delta R_-,\quad f_N^0=\Delta R_-.
\tag{DCA6.3}
\] For \(x\in C\), set \[
z_x=R_-ax,\qquad
w_n(x)=(0,z_x,2z_x,\ldots,(n-1)z_x)\in V_n.
\tag{DCA6.4}
\] All coordinates are retained. Direct componentwise subtraction gives
\[
(\mathsf t-1)w_n(x)=ne_0z_x-\Delta z_x.
\tag{DCA6.5}
\] The zeroth component is \((n-1)z_x\), and each of the other \(n-1\)
components is \(-z_x\). Therefore \[
r_nf_{\mathrm{st}}-f_Nr_1=d_{N_n}w_n+w_nd_{S^{\mathrm{st}}}.
\tag{DCA6.6}
\] In degree \(-1\) both sides vanish; in degree zero the equality is
(DCA6.2) and (DCA6.5). Formula (DCA3.5) now constructs (DCA6.1) from
these continuous maps and this specified homotopy. No pullback of a
singular current at the branch point has been used.

For reference, if a restriction-fibre map has coefficient maps
\(f_S,f_N\) and homotopy \(k\) satisfying \(r'f_S-f_Nr=dk+kd\), its
formula is \[
(s,t)\longmapsto(f_Ss,f_Nt+ks).
\tag{DCA6.7}
\] Substitution into \(d(s,t)=(ds,rs-dt)\) proves the chain identity.
Applying this formula gives \[
\boxed{\mathsf L_{n,\Psi}=\Delta R_-,\qquad
\mathsf L_{n,\Phi}(x)=(L_nx,[w_n(x)]).}
\tag{DCA6.8}
\] The raw-cut correction \([w_n(x)]\) is necessary: omitting it would
give variation \(ne_0R_-ax\) instead of \(\Delta R_-ax\). With it, \[
\mathcal V_n(L_nx,[w_n(x)])
=e_0aL_nx-(\mathsf t-1)w_n(x)=\Delta R_-ax.
\tag{DCA6.9}
\] The map also lands in monodromy invariants, since
\(\mathsf t w_n(x)-e_0aL_nx-w_n(x)=-\Delta z_x\), whose class in \(D_n\)
is zero.

The complete local coefficient maps are \[
\boxed{i^*\mathsf L_n=(R_-,L_n),\qquad
i^!\mathsf L_n=(L_n,nR_-):[C\xrightarrow{a}B]\to[C\xrightarrow{a}B].}
\tag{DCA6.10}
\] The positive angular map is \(nR_-\), because its class in (DCA3.3)
is \(\Sigma_n\Delta R_-=nR_-\). It is consequently also the positive
Dirac support factor. These are the generic/angular/Dirac factors \[
(R_-,\ nR_-,\ nR_-),
\tag{DCA6.11}
\] with their actual domains, rather than three scalar assignments to
one functor.

One can display the invariant copy of \(C\) inside the full vanishing
space without discarding the complement: \[
\iota_n(x)=\left(x,\left[
0,\frac{ax}{n},\frac{2ax}{n},\ldots,\frac{(n-1)ax}{n}
\right]\right).
\tag{DCA6.12}
\] It satisfies \(\mathsf t_\Phi\iota_n=\iota_n\) and
\(\mathcal V_n\iota_n=\Delta a/n\), by the same componentwise
calculation. Formula (DCA6.8) is exactly \(\iota_nL_n\), because
\(aL_n=nR_-a\). All fractions in this comparison are displayed; (DCA6.4)
remains the actual chosen cut correction.

\subsection{DCA7. Both compositions, including their
homotopies}\label{dca7.-both-compositions-including-their-homotopies}

The composite in the target direction is \[
\boxed{\mathsf F_n\mathsf L_n=n\,1_{\mathscr E}}
\quad\text{in the derived sheaf category}.
\tag{DCA7.1}
\] On the stalk complexes its factors are \(nRR_-=n\) and \(S_nL_n=n\).
On the punctured cochain complexes they are \(R\Sigma_n\Delta R_-=n\).
Its gluing homotopy has the remaining value \[
k(x)=R\Sigma_nw_n(x)
=\sum_{j=0}^{n-1}j\,ax
=\frac{n(n-1)}2\,ax.
\tag{DCA7.2}
\] This term has not been suppressed. The degree-minus-one stalk
homotopy \(\ell:C\to B\), \(\ell(x)=\sum_{j=0}^{n-1}j\,ax\), removes it,
because the target restriction in degree \(-1\) is \(r_1^{-1}=1_B\).
More explicitly, for the restriction fibre, \(H(s,t)=(\ell s,0)\) has
\(dH+Hd=(0,ks)\). The same homotopy on the supported summand of the
reconstruction (DCA3.5) proves (DCA7.1) at the sheaf-cone level.

Let \(\mathsf d\) be the positively chosen deck generator on
\(b_{n*}\mathscr E\). It is identity on the stalk model and
\(\mathsf t\) in both punctured degrees. Its gluing homotopy is
\(-e_0a\), because \[
r_n-\mathsf t r_n=-(\mathsf t-1)e_0a.
\tag{DCA7.3}
\] Its induced vanishing map is exactly \(\mathsf t_\Phi\) in (DCA3.8).
The reverse composite is \[
\boxed{\mathsf L_n\mathsf F_n
=\sum_{k=0}^{n-1}\mathsf d^{k}.}
\tag{DCA7.4}
\] On full nearby coefficients both sides are \[
\Delta\Sigma_n=\sum_{k=0}^{n-1}\mathsf t^k:V_n\to V_n.
\tag{DCA7.5}
\] On the pole stalk they are \(n\,1\). On vanishing coefficients the
composite from (DCA4.7) and (DCA6.8) is \[
(x,[v])\longmapsto
\left(nx,\left[(0,ax,2ax,\ldots,(n-1)ax)\right]\right).
\tag{DCA7.6}
\] For the deck norm, sum (DCA3.9). The term
\(\sum_k\mathsf t^kv=\Delta\Sigma_nv\) has zero class in \(D_n\), and
the correction is \[
-\sum_{\ell=0}^{n-2}(n-1-\ell)\mathsf t^\ell e_0ax.
\tag{DCA7.7}
\] Its coordinates are \((-(n-1)ax,-(n-2)ax,\ldots,-ax,0)\). Adding
\((n-1)\Delta ax\) gives exactly the vector in (DCA7.6), proving
equality in the quotient.

At the gluing-homotopy level, the difference between the composite
correction and (DCA7.7) is precisely \((n-1)\Delta ax\). The stalk
homotopy \(\ell(x)=(n-1)ax\) removes it through \(r_n^{-1}=\Delta\),
exactly as in the proof of (DCA7.1). Thus (DCA7.4) holds on the derived
sheaf objects, not merely on their cohomology groups. Likewise
\(\mathsf d^n=1\): its accumulated homotopy \(-\Delta a\) is removed by
the stalk homotopy \(-a\).

For \(n>1\), the deck norm is not generally multiplication by \(n\) on
\(V_n\): it takes \((v_0,\ldots,v_{n-1})\) to \(\Delta(\sum_jv_j)\). The
actual nonzero coefficient space \(B=A'\) is retained in this assertion.
On invariants its value is \(n\), and on the extra cyclic components it
vanishes. On costalks both compositions are \(n\) in each degree,
consistently with (DCA4.7), (DCA6.10), and the identity deck action on
the single supported point.

\subsection{DCA8. The original dual extension and all its comparison
arrows}\label{dca8.-the-original-dual-extension-and-all-its-comparison-arrows}

Classical human-source attribution: extension of continuous linear functionals uses Hahn--Banach; see Terence Tao, \href{https://terrytao.wordpress.com/2009/01/26/245b-notes-6-duality-and-the-hahn-banach-theorem/}{245B Notes 6, Theorem 1 and its complex proof}. A continuous seminorm and its kernel reduce the locally convex application to that normed-space result. This does not assert a continuous splitting.

The closed original source gives the full exact continuous-transpose row
\[
0\longrightarrow Q'\xrightarrow jB\xrightarrow{k'}J'\longrightarrow0,
\qquad k:J\hookrightarrow A.
\tag{DCA8.1}
\] Its surjectivity is continuous Hahn--Banach extension from the closed
subspace \(J\). This does not assert that \(k':B_\beta\to J'_\beta\) is
open.

The original dilation preserves \(J\), so (DCA8.1) commutes with
\((T_n^Q)',T_n',(T_n|_J)'\), and also with their inverse-parameter maps
and the displayed factor \(n\). The actual receiver \[
\mathscr K_r=\operatorname{Cone}(-\mathrm{Gys}_{Q'}\sigma_r)[-1]
\tag{DCA8.2}
\] has the proved triangle \[
\mathscr K_r\longrightarrow\mathscr E_r
\longrightarrow\underline{J'}[1]\longrightarrow\mathscr K_r[1].
\tag{DCA8.3}
\] In continuous current representatives its first map is
\((v,x)\mapsto(q^\vee v,x)\), and its second is
\((b,x)\mapsto k^\vee b\). The coefficient identity \(k'j=0\) proves
their composite zero. Cancelling their shared supported \(C\) coordinate
reduces the cone comparison to the constant-current coefficient row
(DCA8.1), proving the triangle as in RGR9.

Every construction DCA3--DCA7 applies with \(B\) replaced by \(Q'\),
\(a\) by \(\sigma_r\), and with its actual transpose dilation.
Naturality of \(q^\vee,k^\vee\), of the positive Gysin map, and of the
finite sheet sums proves that both cover directions commute with
(DCA8.3). In particular the first map uses \(j\) on each of the \(n\)
nearby sheets and identity on the original \(C\) component; the raw-cut
corrections commute because \(j\sigma_r=a\). The quotient \(J'\) has the
corresponding constant-sheaf trace and unit maps.

There is also the actual map \(\mathscr E_r\to D_cP\), where
\(P=\ker(\underline A\to i_*Q)[1]\). In positive quotient coordinates it
uses \(1_B\) on nearby coefficients and \(\sigma_r\) on vanishing
coefficients. Its cone is \(i_*\operatorname{Cone}(\sigma_r)\), as
proved in RGR5. The cover directions commute with this map by \[
\sigma_rS_n=(T_n^Q)'\sigma_r,\qquad
\sigma_rL_n=n(T_{1/n}^Q)'\sigma_r.
\tag{DCA8.4}
\] Its raw-cut correction commutes componentwise by the same two
identities. Consequently the cover calculation acts on the receiver of
the original dual extension; it has not introduced an unrelated
coefficient extension.

\subsection{DCA9. Both endpoint lines and the actual GTR
factors}\label{dca9.-both-endpoint-lines-and-the-actual-gtr-factors}

The original pole source is \[
V_+=S\oplus\mathbb C_{c_0}\oplus\mathbb C_{c_1},\qquad
V_-=S\oplus\mathbb C_{d_0}\oplus\mathbb C_{d_1},
\] \[
\rho_+(t)(f,c_0,c_1)=(f(\cdot/t),c_0,tc_1),\qquad
\rho_-(t)(h,d_0,d_1)=(t h(t\,\cdot),td_0,d_1).
\tag{DCA9.1}
\] The four labels remain the original value-at-zero and real-integral
labels at their respective poles. Their dual coordinate lines are
denoted by the corresponding primed labels. The full receiver is \[
\mathscr E_{r,\mathrm{full}}=\mathscr E_r\oplus i_*E_p'[-1].
\tag{DCA9.2}
\] The endpoint summand has zero nearby object, stalk and costalk
\(E_p'[-1]\), and shifted vanishing object \(E_p'[-1]\). Its unshifted
vanishing complex is \(E_p'[0]\). Thus the endpoints occur in stalk and
support degree \(1\), separately from the two-term supported complex in
(DCA2.6).

The actual CSD/GTR maps on this summand are \[
\mathsf F_n|_{E_p'}=\rho_p(n)'|_{E_p'},\qquad
\mathsf L_n|_{E_p'}=n\rho_p(1/n)'|_{E_p'}.
\tag{DCA9.3}
\] In the full retained order \((c_0',c_1',d_0',d_1')\), these are \[
\boxed{\mathsf F_n:(1,n,n,1),\qquad
\mathsf L_n:(n,1,1,n).}
\tag{DCA9.4}
\] They follow by transposing the two explicit diagonal endpoint blocks
in (DCA9.1); each product is \(n\). A deck rotation fixes the pole and
acts identically on each endpoint coordinate, so its norm there is
\(n\), also proving the reverse composition on this retained summand.

For comparison, the full original global coefficient complex has \[
D=[V_+\oplus V_-\xrightarrow{r_+-r_-}A\xrightarrow0A]
\quad\text{in degrees }0,1,2,
\] \[
r_+=\Sigma,\qquad r_-=R_A\Sigma=\Sigma\widehat{\phantom f},
\quad R_Ab(u)=u^{-1}b(u^{-1}).
\tag{DCA9.5}
\] Its positive geometric pullback and weighted ramified trace with
inverse coefficients have the exact cochain representatives \[
\mathsf B_n=(\rho(n),T_n,nT_n),\qquad
\mathsf U_n=(n\rho(1/n),nT_{1/n},T_{1/n}).
\tag{DCA9.6}
\] The weighted trace sums all interior sheets and is \(n\) times
identity at each pole, including both endpoint lines. This gives
\(\operatorname{Tr}_n u_n=n1\) on every stalk; its global factors then
follow from the angular and top-degree orientation integrals, as
calculated in GTR2--GTR4.

Transposition of these degree-zero maps adds no Hom-differential sign.
The full original dual rows are \[
\begin{array}{c|ccc}
\text{dual degree}&-2&-1&0\\\hline
\mathsf B_n'&nT_n'&T_n'&\rho(n)'\\
(\mathsf B_n^{-1})'&n^{-1}T_{1/n}'&T_{1/n}'&\rho(1/n)'\\
\mathsf U_n'&T_{1/n}'&nT_{1/n}'&n\rho(1/n)'.
\end{array}
\tag{DCA9.7}
\] All three rows are retained. The second is the pairing contragredient
and differs from the third by the full factor \(n\). The current
pushforward direction (DCA4.3) has the first row's geometry; the
constructed opposite direction (DCA6.1) has the third row's geometry.
Their endpoint maps are exactly (DCA9.3). This establishes the local
cover structure on the new cone. The separate literal dual
identification is supplied by ACD1--ACD4: put
\(b_r=D_r^*\mathsf JEq:A\to H\) and
\(\mathscr P_r=[\underline A^{-1}\xrightarrow{+b_r\mathrm{ev}_p}i_*H^0]\).
The fine resolution ACD2.2 has continuous compact-test dual
chain-isomorphic to \(\mathscr E_r\) by the degree signs \((-1,+1,-1)\)
and the Riesz map \(\overline H\to H'\). Its full pairing is ACD3.3, and
the map \(\mathscr E_r\to D_cP\) is the transpose of the actual
original-source map ACD4.3, with the point-coordinate minus in ACD4.4.
These explicit comparisons identify the primal object and its pairing;
the identification is not inferred from the table of cohomology factors.

\subsection{DCA10. The actual arithmetic image receives matching
actions}\label{dca10.-the-actual-arithmetic-image-receives-matching-actions}

The conclusions about possible separation can be checked directly on the
original arithmetic map rather than on a substitute example. For the
actual value vector \(e_\rho\in H\), the definition (DCA1.7) gives \[
\sigma_r(\overline e_\rho)
=m_\rho\,\overline{d_r(\rho)}\,\delta_{\rho^\#,0},\qquad
a(\overline e_\rho)
=m_\rho\,\overline{d_r(\rho)}\,q'\delta_{\rho^\#,0}.
\tag{DCA10.1}
\] Here \(\delta_{\eta,0}(F)=(\Theta_QF)(\eta)\), so the original Mellin
factor \(1/2\) is part of this functional. Its multiplicity and the full
original defect coefficient have both been retained. The equality
follows by putting \(D_re_\rho=d_r(\rho)e_\rho\) into the weighted
reflected inner product.

The two exact covariance eigenvalue comparisons on this displayed
coordinate are \[
\begin{array}{c|c|c}
\text{map}&\text{source coefficient on }\overline e_\rho
&\text{target coefficient on }q'\delta_{\rho^\#,0}\\\hline
aS_n=Ra&n^{1-\overline\rho}&n^{\rho^\#}\\
aL_n=nR_-a&n^{\overline\rho}&n^{1-\rho^\#}.
\end{array}
\tag{DCA10.2}
\] Each row agrees identically because \(\rho^\#=1-\overline\rho\). The
coefficients are equal whether or not \(d_r(\rho)\) vanishes. When it
vanishes, (DCA10.1) is the zero map at that actual coordinate; when it
does not, (DCA10.2) is exact equivariance on its nonzero image line. No
existence of an off-critical zero is assumed.

The closed source kernel is \[
\ker a=\overline{\ker D_r}
=\{\overline y:y_\rho=0\text{ whenever }\Re\rho\ne1/2\}.
\tag{DCA10.3}
\] The equality follows from injectivity of \(j,A_H\), and from
\(r^{\Re\rho}=r^{1-\Re\rho}\) if and only if \(\Re\rho=1/2\). Both
actions in (DCA1.13) preserve this exact kernel and the actual image
\(a(C)\). Their induced quotient maps are continuous for quotient
topologies, without any conclusion that the image is closed.

Thus the discrepancy between the constant-current factor \(nR\) and the
angular/support factor \(R\) does not impose a zero attaching map. Those
factors are the two receivers (DCA5.1)--(DCA5.2) of the full nearby map
\(R\Sigma_n\). The supported arithmetic map has precisely the matching
eigenvalues in (DCA10.2). To demand an incompatible weight on that same
supported map would alter its already proved action. In particular, this
cover calculation has not proved \(a=0\) or \(D_r=0\).

\subsection{DCA11. The exact continuation supplied by the
calculation}\label{dca11.-the-exact-continuation-supplied-by-the-calculation}

The resulting operation on the actual Gysin obstruction
\(\mathfrak o_r=\mathrm{Gys}_Ba\) is now computed, together with its
full nearby source, both vanishing maps, costalk maps, endpoint maps,
and the connecting morphisms to the dual original extension. Under
support adjunction it obeys \(Ra=aS_n\) in the proper-current direction
and \(nR_-a=aL_n\) in the opposite direction. These are identities on
the complete source \(\overline H\), not merely on the displayed finite
value coordinates.

The next receiver after the automatically vanishing invariant-cycle
image quotient is the actual image \(a(C)\), as computed in ADC3. Its
source is the exact quotient \(C/\ker a\), and DCA10 gives its existing
equivariant map into \(q'Q'\subset A'\). The concrete cover attempt is
complete: its action does not separate this source from this target,
because the actions match exactly. The remaining calculation must use
additional proved structure of that same image or of its original
continuous pairing, while retaining (DCA1.13), (DCA8.3), and all
endpoint terms. Replacing the supported target factor by its invariant
factor would contradict (DCA5.3), rather than supply such a calculation.

\subsection{DCA12. Scope and reproducible algebraic
checks}\label{dca12.-scope-and-reproducible-algebraic-checks}

Every finite-sheet identity used in the proof can be checked by the
following displayed component calculations: \((\mathsf t-1)\Delta=0\);
\(\Sigma_n(\mathsf t-1)=0\); the partial-sum inverse in DCA3;
\((\mathsf t-1)(0,z,\ldots,(n-1)z)=ne_0z-\Delta z\);
\(\sum_k\mathsf t^k=\Delta\Sigma_n\); and the coordinate difference
\((n-1)\Delta ax\) between (DCA7.6) and (DCA7.7). They prove the signs,
the retained raw-cut correction, and both derived composition laws for
every positive integer \(n\), including \(n=1\), when the extra sheet
quotient \(D_n=B^n/\Delta B\) is zero, the correction vector is zero,
and both cover maps are identity. The independently fixed arithmetic
parameter remains \(r>1\).

The argument constructs all stated sheaf maps from continuous currents
or explicit continuous restriction-cone data. It never infers a unique
current chain map from cohomology alone. No conclusion about strong
openness of \(A'_\beta\to J'_\beta\), image closure of \(a(C)\),
finite-dimensionality, arithmetic purity, or vanishing of the actual
adjoint defect is inserted. Every original coefficient and endpoint
contribution used in the construction remains present in the formulas
above.
